% !TeX root = ../../Book2.tex
% 纲要标题：### C.1 中心单代数的交换基环推广
% 纲要位置：工作记录/计划纲要.md，第 1826 行。

% 纲要要点（供目录核对）：
% * 用忠实有限生成投射模与左右乘包络映射推广中心单代数；分清单环、Azumaya、原基环分裂及矩阵代数
% #### C.1.1 Azumaya 代数的定义
% #### C.1.2 域情形与矩阵代数

\section{中心单代数的交换基环推广}
\label{b2:sec:C01}

中心单代数的定义包含域上的有限维性，但若把底域换成交换环，有限生成模未必有基，代数本身也通常不再是单环。例如 \(M_n(\mathbb Z)\) 有非零真双边理想 \(2M_n(\mathbb Z)\)，但在每个素数 \(p\) 处取剩余域纤维，仍得到中心单代数 \(M_n(\mathbb F_p)\)。所以推广的对象不能要求自身是单环，而要让各点上的中心单结构来自同一个有限投射模，并由左右乘作用统一描述。本附录的基环 \(R\) 交换、含幺且非零，代数的结构同态取值于中心；底环变换的目标环也约定非零。

\subsection{Azumaya 代数的定义}
\label{b2:subsec:C01:01}

对一个 \(R\) 代数 \(A\)，其反代数 \(A^{\mathrm{op}}\) 使用相反次序的乘法。左右乘给出 \(R\) 代数同态
\begin{equation}\label{b2:eq:azumaya-enveloping}
\mu_A:A\otimes_R A^{\mathrm{op}}\longrightarrow\operatorname{End}_R(A),
\qquad a\otimes b^{\mathrm{op}}\longmapsto(x\mapsto axb).
\end{equation}
例如，两次作用的复合把 \(x\) 送到 \(aa'xb'b\)，这正对应张量代数中的乘积 \(aa'\otimes(b'b)^{\mathrm{op}}\)。反代数保证右乘一侧的次序正确。

\begin{definition}\label{b2:def:azumaya-algebra}
设 \(R\) 为非零交换含幺环、\(A\) 为结合含幺 \(R\) 代数，其结构同态的像位于 \(A\) 的中心。若 \(A\) 满足以下两项条件，则称它为 \(R\) 上的\term[Azumaya-daishu]{Azumaya 代数}[Azumaya algebra]：
\begin{enumerate}
\item 作为 \(R\) 模，\(A\) 忠实、有限生成且投射。这里忠实指 \(rA=0\) 蕴含 \(r=0\)；有限生成且投射等价于存在一个 \(R\) 模 \(Q\) 及正整数 \(N\)，使 \(A\oplus Q\cong R^N\)。
\item 左右乘给出的 \(R\) 代数同态 \(\mu_A:A\otimes_RA^{\mathrm{op}}\to\operatorname{End}_R(A)\)、\(a\otimes b^{\mathrm{op}}\mapsto(x\mapsto axb)\) 是同构。
\end{enumerate}
\end{definition}
\symidx{\(\mu_A\)：Azumaya 代数的左右乘作用映射}

域上\cref{b2:thm:csa-enveloping-isomorphism}正是用同一个包络映射刻画中心单代数；这里将底域换成交换环，并用忠实有限生成投射模代替非零有限维向量空间。定义中的直和刻画由\cref{b2:prop:finite-projective-summand}保证。它允许模没有全局基，却仍能在局部选择有限基。忠实性则使结构同态 \(R\rightarrow A\) 单射：若 \(r1_A=0\)，就有 \(rA=0\)。因此可以把 \(R\) 看作 \(A\) 中的标量。\cref{b2:prop:azumaya-center}将进一步证明，\(A\) 的中心恰好是这些标量。

\ProseParagraphBreak
\(R\) 本身是 \(R\) 上的 Azumaya 代数，因为 \(R\otimes_R R\cong R\cong\operatorname{End}_R(R)\)，且这些同构把 \(\mu_R\) 变成恒等映射。若 \(0\ne I\subsetneq R\)，商环 \(R/I\) 的标量作用却零化了整个 \(I\)，不能作为同一个基环 \(R\) 上的 Azumaya 代数。例如取 \(R=k\times k\)、\(I=0\times k\)，则 \(A=R/I\cong k\) 是有限投射 \(R\) 模，包络映射也同构，却完全遗漏了基环的第二个分量。这说明忠实性是独立需要的条件。改变底环后，\(R/I\) 作为自身上的代数满足定义；底环是概念的一部分。

\subsection{域情形与矩阵代数}
\label{b2:subsec:C01:02}

\begin{proposition}\label{b2:prop:azumaya-field-matrix}
域 \(k\) 上的 Azumaya 代数恰为有限维中心单 \(k\) 代数。对任意非零交换环 \(R\) 及 \(n\ge1\)，\(M_n(R)\) 都是 \(R\) 上的 Azumaya 代数。
\end{proposition}
\begin{proof}
域上的忠实有限生成投射模恰为非零有限维向量空间。在此前提下，\cref{b2:thm:csa-enveloping-isomorphism}说明，左右乘作用同构等价于代数中心单。这给出第一个断言的两个方向。

对 \(A=M_n(R)\)，矩阵单位 \(E_{ij}\) 给出 \(A\) 的 \(R\) 基，所以它忠实且有限自由。直接计算得
\[
\mu_A(E_{ij}\otimes E_{\ell m}^{\mathrm{op}})(E_{pq})
=\delta_{jp}\delta_{q\ell}E_{im}.
\]
随着 \(i,j,\ell,m\) 变化，右侧恰为自由模 \(A\cong R^{n^2}\) 的自同态环中的标准矩阵单位：它把一个指定基向量 \(E_{j\ell}\) 送到一个指定基向量 \(E_{im}\)，并把其余基向量送到零。因此 \(\mu_A\) 把来源的一组基双射到目标的一组基，故为同构。
\end{proof}

例如 \(M_2(\mathbb Z/4\mathbb Z)\) 同样满足定义：基环可以有零因子，矩阵大小也不必在基环中可逆。Azumaya 性不是把矩阵迹除以矩阵大小所得到的性质。

\ProseParagraphBreak
另一端的例子是实四元数代数 \(\mathbb H\)。\cref{b2:prop:quaternion-division-or-split}及范数公式\eqref{b2:eq:quaternion-norm}说明，它是中心除 \(\mathbb R\) 代数：非零四元数的范数是四个实数平方之和，因而不为零。所以 \(\mathbb H\) 是 \(\mathbb R\) 上的 Azumaya 代数，却不同构于 \(M_2(\mathbb R)\)，因为后者有非零零因子。定义允许这种在原基域上未分裂的中心单代数；要求 Azumaya 性，并没有要求原基环上已经存在矩阵单位。
